% DF-001, CF-007/012, P09 del inventario REV05.
% Residencia: APP toroidal y transporte de cursor; ejemplo tras el cociclo.
% Propietarios: PAQUETE_ARTICULO_I_INTEGRACION_ESTADO_CAMPOS_20260919/
% deltas/integracion/APPRouteWinding.lean (íntegro),
% lean/biblioteca/TPKTransport.lean:19--143, 155--229,
% lean/biblioteca/TPKFiniteCursor.lean:17--146.
% Estatuto: RESULTADO_RECUPERADO; desarrollo del ejemplo finito efectivo,
% sin identificar el factor observable con todo el estado enriquecido.

\section[Retour spatial et enroulement du curseur]{Retour spatial et conservation de l'enroulement dans le calendrier de vingt-sept transitions}
\label{rc27:sec:cursor}

Les deux coordonnées périodiques de l'APP proviennent de deux coordonnées
entières transportées. La lecture résiduelle identifie les positions qui
diffèrent de multiples de neuf ; les quotients entiers conservent les
tours de ce transport. Le calendrier TPK fournit un exemple
complet dans lequel la position résiduelle revient et le curseur relevé
conserve un déplacement horizontal non nul.

\subsection{Coordonnées relevées et cocycle de frontière}

Pour $z\in\mathbb Z$, soit
\[
\rho_9(z)=1+((z-1)\bmod9),\qquad
q_9(z)=\left\lfloor\frac{z-1}{9}\right\rfloor.
\]
La division euclidienne, également pour les coordonnées négatives, donne
\begin{equation}
z=\rho_9(z)+9q_9(z),\qquad1\le\rho_9(z)\le9.
\label{rc27:eq:division}
\end{equation}
La retenue de transport d'un incrément $a\in\mathbb Z$ est
\[
c(z,a)=q_9(z+a)-q_9(z).
\]
De \eqref{rc27:eq:division} se déduisent
\begin{align}
\rho_9(z+a)+9c(z,a)&=\rho_9(z)+a,\\
c(z,a+b)&=c(z,a)+c(z+a,b).
\label{rc27:eq:cociclo}
\end{align}
La seconde identité découle de l'annulation de $q_9(z+a)$ dans le membre
de droite. Le second incrément est évalué depuis l'extrémité du premier ;
cette dépendance conserve l'information du franchissement de frontière.

Soit $D=\{N,E,S,O\}$, avec les déplacements
\[
\delta(N)=(-1,0),\quad\delta(E)=(0,1),\quad
\delta(S)=(1,0),\quad\delta(O)=(0,-1).
\]
Pour un mot de directions $w=d_1\cdots d_n$, définissons
$\Delta(w)=\sum_j\delta(d_j)$ et, depuis $o=(o_1,o_2)$,
\[
\operatorname{Ext}(o,w)=o+\Delta(w),\qquad
\nu(o,w)=\bigl(c(o_1,\Delta_1(w)),c(o_2,\Delta_2(w))\bigr).
\]
Le vecteur $\nu$ enregistre les deux enroulements.

\begin{proposition}[Composition et inversion des tours]
Pour deux routes concaténables $u,v$, on a
\[
\nu(o,uv)=\nu(o,u)+\nu(\operatorname{Ext}(o,u),v).
\]
Si $w^{-1}$ inverse l'ordre et remplace chaque direction par son
opposée, alors
\[
\operatorname{Ext}(\operatorname{Ext}(o,w),w^{-1})=o,
\qquad
\nu(\operatorname{Ext}(o,w),w^{-1})=-\nu(o,w).
\]
Pour une route spatialement fermée dans l'APP,
\[
\nu(o,w)=\Delta(w)/9\in\mathbb Z^2.
\]
\end{proposition}
\begin{proof}
La somme des déplacements est additive sous concaténation. Appliquer
\eqref{rc27:eq:cociclo} dans chaque coordonnée démontre la première
identité. L'inversion donne $\Delta(w^{-1})=-\Delta(w)$ et récupère
l'extrémité initiale. Soustraire les quotients dans l'ordre contraire change
le signe de $\nu$. Enfin, la fermeture spatiale équivaut à
$\rho_9(o_i+\Delta_i)=\rho_9(o_i)$ pour les deux coordonnées ; la première
égalité de \eqref{rc27:eq:cociclo} donne alors $9\nu_i=\Delta_i$.
\end{proof}

Les cycles de neuf mouvements vers le sud et de neuf mouvements vers l'est,
depuis $(1,1)$, ont pour tours $(1,0)$ et $(0,1)$, respectivement.
Ce sont deux cycles indépendants de la surface périodique APP.

\subsection{Mise à jour des deux curseurs}

Un curseur observable relevé est $c=(x,y,d)\in\mathbb Z^2\times D$.
La phase $p$ appartient à $\mathbb Z/108\mathbb Z$ et est représentée
par $0,\ldots,107$. Dans le feuillet additif, le déplacement est activé
lorsque $p\equiv0\pmod3$ ; dans le multiplicatif, lorsque
$p\equiv1\pmod3$. Définissons
\[
a_+(p)=\mathbf1_{\{p\equiv0\ (3)\}},\qquad
a_\times(p)=\mathbf1_{\{p\equiv1\ (3)\}},\qquad
\varepsilon(p)=\mathbf1_{\{p+1\equiv0\ (27)\}}.
\]
L'opérateur observable agit au moyen de
\begin{equation}
c'_s=\bigl(x_s+a_s(p)\delta_1(d_s),\,
           y_s+a_s(p)\delta_2(d_s),\,
           \operatorname{op}^{\varepsilon(p)}d_s\bigr),
\qquad p'=p+1\pmod{108},
\label{rc27:eq:actualizacion}
\end{equation}
où $s\in\{+,\times\}$ et $\operatorname{op}$ échange
$N$ avec $S$ et $E$ avec $O$.
Le déplacement utilise la direction précédente, puis la
transition de demi-cycle met à jour l'orientation.

L'inverse récupère d'abord $p=p'-1\pmod{108}$, restitue
$d_s=\operatorname{op}^{\varepsilon(p)}d'_s$ et soustrait
$a_s(p)\delta(d_s)$ des coordonnées. Comme
$\operatorname{op}^2=I$, les deux compositions sont l'identité.

\begin{theorem}[Retour spatial effectif dans la première demi-couronne]
\label{rc27:thm:ejemplo}
Partons de
$c_+(0)=c_\times(0)=(1,1,E)$ et $p(0)=0$.
Après vingt-sept transitions de
\eqref{rc27:eq:actualizacion}, on obtient
\[
c_+(27)=c_\times(27)=(1,10,O),\qquad p(27)=27.
\]
Les deux positions résiduelles reviennent à $(1,1)$, tandis que le
quotient horizontal augmente d'une unité et que le quotient vertical demeure
invariable.
\end{theorem}
\begin{proof}
Pour $0\le n\le27$, le nombre d'activations additives antérieures
à la transition $n$ est $\lfloor(n+2)/3\rfloor$ ; le nombre
d'activations multiplicatives est $\lfloor(n+1)/3\rfloor$. Avant
la transition de phase de $26$ à $27$, il n'y a pas d'inversion de direction.
Par récurrence sur $n$, les coordonnées sont
\[
x_+(n)=x_\times(n)=1,\qquad
y_+(n)=1+\left\lfloor\frac{n+2}{3}\right\rfloor,\qquad
y_\times(n)=1+\left\lfloor\frac{n+1}{3}\right\rfloor.
\]
Les directions sont $E$ pour $0\le n<27$ et $O$ pour $n=27$.
La décomposition dans chaque groupe de trois transitions est
\[
\begin{array}{c|c|c|c}
n&y_+(n)&y_\times(n)&\text{direction}\\\hline
3k&1+k&1+k&E\\
3k+1&2+k&1+k&E\\
3k+2&2+k&2+k&E
\end{array}
\quad(0\le k\le8),
\]
et l'extrémité finale est $(y_+,y_\times,d)=(10,10,O)$.
Il y a neuf activations dans chaque feuillet. La transition finale a pour phase
préalable $26$, de classe deux modulo trois ; elle inverse donc les
directions et maintient les coordonnées qui ont déjà atteint dix.

Enfin, $\rho_9(10)=1$, $q_9(10)=1$,
$\rho_9(1)=1$ et $q_9(1)=0$. La position résiduelle coïncide avec
la position initiale, et l'enroulement spatial est $(0,1)$.
\end{proof}

Le retour de cet exemple est spatial. La direction et la phase
enregistrées sont $(O,27)$ à la fin et $(E,0)$ au début. L'exemple
conserve expressément les trois données : position résiduelle, orientation et
phase ; leur égalité ou leur différence est établie pour chacune.

\subsection{Récupération du curseur et compatibilité à toute profondeur}

\begin{proposition}[Récupération du curseur relevé]
La position résiduelle, la direction et les deux quotients entiers
déterminent de manière unique le curseur $c=(x,y,d)$.
\end{proposition}
\begin{proof}
Les deux identités \eqref{rc27:eq:division} récupèrent $x$ et $y$
et la direction est conservée dans la projection observable. Si ces
cinq données coïncident pour deux curseurs, leurs trois
composantes coïncident et, par conséquent, les curseurs sont égaux.
\end{proof}

Soit $P$ la projection des deux curseurs sur leurs positions modulo
neuf, en conservant leurs directions et la phase. La mise à jour
finie $T_{\rm fin}$ utilise les mêmes indicateurs d'activation et
d'orientation, avec une somme résiduelle dans les coordonnées. Si $T$ est
\eqref{rc27:eq:actualizacion}, alors
\[
PT=T_{\rm fin}P,\qquad PT^n=T_{\rm fin}^{\,n}P
\quad(n\ge0).
\]
La première égalité résulte de la compatibilité de la réduction
modulo neuf avec l'addition et de la conservation de la direction et de la phase ;
la seconde se démontre par récurrence sur $n$.

L'archive ordonnée des états observables peut être prolongée au moyen de
$(q,\mathcal A)\mapsto(Tq,\mathcal A\Vert q)$.
Après $n$ transitions, elle conserve son préfixe initial et sa longueur
augmente de $n$. Cette opération maintient séparés l'état courant,
l'enroulement et le registre de ses états antérieurs. Dans la
construction enrichie complète, le facteur observable se compose
avec les autres lectures de cylindre, d'incidence et de mémoire au moyen de leurs
projections déclarées.
